
4) 2. Caso general.\par\nopagebreak
Se ve, utilizando el lema~\Ref{XIII.6.1.1}, 4)\kern2pt1 y
\ifthenelse{\boolean{orig}}
{1)\kern2pt a),}
{1)\kern2pt b),}
que, tras restringir $S$ a un abierto no vacío, puede suponerse que
$S(f_*\Phi)$ es constructible y que $(S\Phi,f)$ es cohomológicamente
propio relativamente a~$S$ en dimensión $\leqslant 0$. Puede entonces
encontrarse un haz representable por un $Y$-esquema étale de tipo
finito~$T$ y un epimorfismo
$$
a \colon T \to S(f_*\Phi)\quoi.
$$
Se retoman las notaciones de 4)\kern2pt1 y se pone además
$Z=T\times_Y T$, $X_Z=X\times_Y Z$; se denota por
$f_Z \colon X_Z \to Z$ el morfismo canónico. Puede suponerse que~$T$
se ha elegido de modo que la imagen~$q$ de la sección identidad de~$T(T)$
por~$a$ esté definida por un objeto~$p$ de
$(f_*\Phi)_T=\Phi_{X_T}$. Sean $p_1$ y $p_2$ (\resp $q_1$ y $q_2$)
las imágenes inversas de~$p$ (\resp $q$) por las dos proyecciones
de~$Z$ en~$T$. Tras restringir $S$ a un abierto no vacío, puede
suponerse que $(\SheafAut_{X_T}(p),f_T)$ es cohomológicamente propio
relativamente a~$S$ en dimensión $\leqslant 1$, que
$(\SheafHom_{X_Z}(p_1,p_2),f_Z)$ es cohomológicamente propio
relativamente a~$S$ en dimensión $\leqslant 0$, y que
$f_{T*}(\SheafAut_{X_T}(p))$ y
$f_{Z*}(\SheafHom_{X_Z}(p_1,p_2))$ son constructibles
\ifthenelse{\boolean{orig}}
{(a)\kern2pt1) y 4)\kern2pt1)).}
{(1)\kern2pt b) y 4)\kern2pt1)).}

Se deduce en primer lugar que, para todo esquema $Y_1$ étale sobre~$Y$
y para todo par de objetos $y$, $y_1$ de~$(f_*\Phi)_{Y_1}$, el haz
$\SheafAut_{Y_1}(y)$ (\resp $\SheafHom_{Y_1}(y,y_1)$) es
constructible, pues tal haz es, localmente para la topología étale
de~$Y_1$, imagen inversa de~$f_{T*}(\SheafAut_{X_T}(p))$
(\resp de~$f_{Z*}(\SheafHom_{X_Z}(p_1,p_2))$).

Queda por demostrar que $(\Phi,f)$ es cohomológicamente propio
relativamente
\marginpar{405}
a~$S$ en dimensión $\leqslant 1$. Para ello basta demostrar que,
\ifthenelse{\boolean{orig}}
{para todo $S'$-esquema~$S$ y para todo punto geométrico
$\overline{y}'$ de~$Y$, si se denota por~$\overline{Y}$ la localización
estricta de~$Y$ en~$y'$, $\overline{X}=X\times_Y\overline{Y}$, etc.,}
{para todo $S$-esquema~$S'$ y para todo punto geométrico
$\overline{y}'$ de~$Y'$, si $\overline{y}$ denota su imagen en~$Y$ y
$\overline{Y}$ (\resp $\overline{Y'}$) denota la localización estricta
de~$Y$ en~$\overline{y}$ (\resp de~$Y'$ en~$\overline{y}'$), y se pone
$\overline{X}=X\times_Y\overline{Y}$,
$\overline{X'}=X'\times_{Y'}\overline{Y'}$, etc.,}
entonces el funtor canónico
$$
\overline{\varphi} \colon \Phi(\overline{X}) \to
\Phi'(\overline{X'})
$$
es una equivalencia de categorías.

Demostremos que~$\overline{\varphi}$ es plenamente fiel. Sean
$x$, $y \in \Phi(\overline{X})$, y sean $x'$, $y'$ sus imágenes en
$\Phi'(\overline{X'})$; demostremos que el morfismo canónico
\begin{equation*}
\label{eq:XIII.4.*} \tag{$*$} {\Hom_{\overline{X}}(x,y)
\to \Hom_{\overline{X'}}(x',y') }
\end{equation*}
es biyectivo. Por definición de~$T$, existen dos morfismos de
$\overline{Y}$ en~$T$ tales que $x$ e~$y$ son las imágenes inversas
de~$p$ por esos dos morfismos. Esto equivale a decir que existe un
morfismo $\overline{Y} \to Z$ y, por tanto, un morfismo
$h\colon \overline{X} \to X_Z$ tal que
$$
h^*(p_1)=x \qquad h^*(p_2)=y.
$$
Por consiguiente, se tiene un isomorfismo canónico
$$
\SheafHom_{\overline{X}}(x,y)=h^*(\SheafHom_{X_Z}(p_1,p_2) )\quoi.
$$
Pero, teniendo en cuenta que $(\SheafHom_{X_Z}(p_1,p_2),f_Z)$ es
cohomológicamente propio relativamente a~$S$ en dimensión
$\leqslant 0$, se ve que lo mismo vale para
$(\SheafHom_{\overline{X}}(x,y),\overline{f})$, lo que demuestra que
el morfismo~\eqref{eq:XIII.4.*} es biyectivo.

Demostremos que~$\overline{\varphi}$ es esencialmente sobreyectivo. Sea
$x'\in \Phi'(\overline{X'})$. Como $(S\Phi,f)$ es cohomológicamente
propio relativamente a~$S$ en dimensión $\leqslant 0$, el morfismo
canónico
$$
\H^0(\overline{X},\overline{S\Phi}) \to
\H^0(\overline{X'},\overline{S\Phi}')
$$
es biyectivo. Sea $G'$ el subgerbe maximal de~$\Phi'$ generado
por~$x'$; existe entonces un subgerbe maximal~$G$ de
$\overline{\Phi}$ cuya imagen inversa sobre~$\overline{X'}$ es~$G'$.
Según 4)\kern2pt1, $(G,\overline{f})$ es cohomológicamente propio
relativamente a~$S$ en dimensión $\leqslant 1$; por consiguiente, el
funtor canónico
$$
G(\overline{X}) \to G'(\overline{X'})
$$
\marginpar{406}%
es una equivalencia de categorías, lo que demuestra la existencia
de un objeto~$x$ de~$\Phi(\overline{X})$ cuya imagen en
\ifthenelse{\boolean{orig}}
{$\Phi'(X')$}
{$\Phi'(\overline{X'})$}
es isomorfa a~$x'$ y concluye la demostración del teorema.
